\documentclass[11pt]{article}
\usepackage[a4paper,margin=1in]{geometry}
\usepackage{array}
\usepackage{booktabs}
\usepackage{longtable}
\usepackage{hyperref}
\usepackage{enumitem}
\usepackage{parskip}

\hypersetup{
  colorlinks=true,
  linkcolor=black,
  urlcolor=blue
}

\setlength{\LTpre}{0pt}
\setlength{\LTpost}{0pt}
\renewcommand{\arraystretch}{1.2}

\begin{document}

\begin{center}
{\LARGE MCP Server Attack Benchmarks and Datasets} \\
\vspace{0.4em}
{\large Focused shortlist from \texttt{reviews/benchmark\_md}} \\
\vspace{0.4em}
Generated for thesis planning and experiment design
\end{center}

\section*{Purpose}
This report extracts the benchmarks and datasets that are most useful when the research question is:
\textbf{how to detect, score, or defend against attacks that reach or abuse the MCP server}.

The shortlist favors resources that do at least one of the following:
\begin{itemize}[leftmargin=1.2em]
\item include direct attacks on MCP servers or server-side attack surfaces,
\item show how malicious tool metadata or protocol behavior turns an agent into a threat to the server,
\item provide server-level prevalence or capability data that helps calibrate risk scoring.
\end{itemize}

\section*{Quick Recommendation}
\begin{longtable}{p{0.22\textwidth} p{0.12\textwidth} p{0.14\textwidth} p{0.42\textwidth}}
\toprule
\textbf{Resource} & \textbf{Type} & \textbf{Priority} & \textbf{Why it matters for MCP server attacks} \\
\midrule
MCPSecBench & Benchmark & Highest & Explicit attack-surface framing, including server-side and protocol-side attacks, with empirical attack success rates. \\
MCP-AttackBench & Dataset & Highest & Largest MCP attack corpus; strong for training classifiers that flag malicious requests before they hit the server. \\
Component-based Attack PoC & Dataset & Highest & Directly models malicious MCP servers and attack components; useful for server-defense and compositional risk. \\
MCPTox & Benchmark & High & Real-world poisoned tool attacks on 45 live MCP servers; strong ecological validity. \\
MCP-SafetyBench & Benchmark & High & Includes server-side attacks across five real-world domains and 13 models in multi-turn settings. \\
MCP Server Database & Dataset & High & Server and tool capability census with exploit-risk combinations; useful for exfiltration and multi-tool escalation logic. \\
MCP Server Dataset (67K) & Dataset & Medium-High & Ecosystem-scale prevalence of server hijacking, shadowing, confusion, and credential leakage. \\
MCP-ITP & Dataset & Medium & Strong for subtle poisoning that causes the agent to behave dangerously toward the server, but indirect rather than server-side. \\
\bottomrule
\end{longtable}

\section*{Tier 1: Best Direct Fits}

\subsection*{1. MCPSecBench}
\textbf{File:} \texttt{05\_mcpsecbench.md} \\
\textbf{Why it is top priority:}
\begin{itemize}[leftmargin=1.2em]
\item It explicitly separates attacks by \textbf{server-side}, \textbf{protocol-side}, \textbf{client-side}, and \textbf{host-side} surfaces.
\item It reports empirical success rates, including about \textbf{47\%} ASR for server-side attacks and \textbf{100\%} ASR for protocol-side attacks.
\item That makes it the cleanest benchmark for assigning severity to attacks that can reach or compromise the server boundary.
\end{itemize}
\textbf{Best use in your project:}
\begin{itemize}[leftmargin=1.2em]
\item Define the attack-surface or server-asset part of your risk score.
\item Justify maximum-severity handling for protocol-level attacks.
\item Evaluate whether a scorer distinguishes direct server threats from less severe layers.
\end{itemize}

\subsection*{2. MCP-AttackBench}
\textbf{File:} \texttt{03\_mcp\_attackbench.md} \\
\textbf{Why it is top priority:}
\begin{itemize}[leftmargin=1.2em]
\item It is the largest MCP security dataset in the folder, with \textbf{70,448} samples.
\item The taxonomy includes \textbf{protocol-specific} and \textbf{injection/execution} attacks, which are highly relevant to server defense.
\item It is directly suitable for training a classifier that decides whether an incoming MCP request should be treated as benign or adversarial.
\end{itemize}
\textbf{Best use in your project:}
\begin{itemize}[leftmargin=1.2em]
\item Train the core request-risk detector.
\item Map attack families into severity bands for your scoring pipeline.
\item Use cascade stage or detection-difficulty ideas to justify fast-path versus deep-analysis scoring.
\end{itemize}

\subsection*{3. Component-based Attack PoC Dataset}
\textbf{File:} \texttt{08\_component\_attack\_poc.md} \\
\textbf{Why it is top priority:}
\begin{itemize}[leftmargin=1.2em]
\item It is directly about \textbf{when MCP servers attack}.
\item The dataset decomposes malicious servers into launch commands, initialization snippets, tools, resources, and prompts.
\item It supports a compositional view of risk: the server can become dangerous through combinations of seemingly separate components.
\end{itemize}
\textbf{Best use in your project:}
\begin{itemize}[leftmargin=1.2em]
\item Build or justify a server-defense model that scores risk from multiple components together.
\item Support arguments about why binary scanners miss real server attacks.
\item Motivate temporal or session-level aggregation in your framework.
\end{itemize}

\subsection*{4. MCPTox}
\textbf{File:} \texttt{04\_mcptox.md} \\
\textbf{Why it is high priority:}
\begin{itemize}[leftmargin=1.2em]
\item It uses \textbf{45 live real-world MCP servers} and \textbf{353 authentic tools}.
\item The attack mechanism is poisoned tool metadata, which is often how malicious intent reaches the agent and then causes risky actions against the server.
\item It is the strongest realism benchmark for tool poisoning on real MCP infrastructure.
\end{itemize}
\textbf{Best use in your project:}
\begin{itemize}[leftmargin=1.2em]
\item Evaluate the \textbf{agent compromise indicator} or similar dimension.
\item Show that attacks on the tool layer can become server risk once the manipulated agent starts issuing dangerous calls.
\item Add real-world credibility to your experimental section.
\end{itemize}

\subsection*{5. MCP-SafetyBench}
\textbf{File:} \texttt{06\_mcp\_safetybench.md} \\
\textbf{Why it is high priority:}
\begin{itemize}[leftmargin=1.2em]
\item It includes \textbf{server-side}, \textbf{host-side}, and \textbf{user-side} attack surfaces across five domains.
\item The multi-turn setup is useful when attacks escalate over a session rather than in one isolated request.
\item It provides cross-model evidence using 13 models, which helps explain why the same server may face different risk depending on the agent.
\end{itemize}
\textbf{Best use in your project:}
\begin{itemize}[leftmargin=1.2em]
\item Evaluate multi-turn or session-aware server risk.
\item Add domain weighting, for example finance versus repository management.
\item Compare model-aware modifiers in your risk-scoring design.
\end{itemize}

\section*{Tier 2: Strong Supporting Datasets}

\subsection*{6. MCP Server Database}
\textbf{File:} \texttt{07\_mcp\_server\_database.md} \\
\textbf{Why it helps:}
\begin{itemize}[leftmargin=1.2em]
\item It classifies tools into \textbf{EIT}, \textbf{PAT}, and \textbf{NAT}, which is useful for server exposure modeling.
\item It shows that \textbf{27.2\%} of servers expose exploitable tool combinations.
\item It is very useful for data exfiltration and cross-tool escalation, even if it is not a benchmark in the narrow evaluation sense.
\end{itemize}
\textbf{Best use in your project:}
\begin{itemize}[leftmargin=1.2em]
\item Build server-side capability priors.
\item Calibrate exfiltration risk and multi-tool attack chains.
\item Support the claim that combinations of tools are often riskier than any single tool.
\end{itemize}

\subsection*{7. MCP Server Dataset (67K)}
\textbf{File:} \texttt{09\_mcp\_server\_dataset\_67k.md} \\
\textbf{Why it helps:}
\begin{itemize}[leftmargin=1.2em]
\item It gives ecosystem-scale prevalence: \textbf{67,057} servers and \textbf{44,499} analyzed Python tools.
\item It documents concrete server-relevant problems such as tool shadowing, tool confusion, credential leakage, and \textbf{111+ server hijacking instances}.
\item It is ideal for grounding your scoring thresholds in real base rates rather than only benchmark results.
\end{itemize}
\textbf{Best use in your project:}
\begin{itemize}[leftmargin=1.2em]
\item Set realistic priors for registry trust and ecosystem risk.
\item Justify why certain tool patterns should start with elevated suspicion.
\item Support the thesis motivation section with prevalence evidence.
\end{itemize}

\subsection*{8. MCP-ITP Implicit Poisoning}
\textbf{File:} \texttt{10\_mcp\_itp\_implicit\_poisoning.md} \\
\textbf{Why it helps:}
\begin{itemize}[leftmargin=1.2em]
\item It focuses on subtle poisoning that is almost invisible to current defenses.
\item That poisoning does not attack the server directly, but it can make the agent behave dangerously toward the server.
\item The near-zero detector success reported in the note makes it important for the ``compromised agent becomes server threat'' argument.
\end{itemize}
\textbf{Best use in your project:}
\begin{itemize}[leftmargin=1.2em]
\item Support the compromise or hostile-influence dimension in the score.
\item Explain why server defense cannot rely only on explicit attack signatures.
\item Evaluate text-based poisoning detection around tool descriptions.
\end{itemize}

\section*{Internet-Sourced Emerging Additions}
The following resources were verified online on \textbf{2026-03-30} and added after the original local benchmark sweep.
They are useful, but they are not as mature or as directly usable as the core Tier 1 resources.

\subsection*{9. MCPShield}
\textbf{File:} \texttt{19\_mcpshield.md} \\
\textbf{Source:} \url{https://arxiv.org/abs/2602.14281}
\begin{itemize}[leftmargin=1.2em]
\item Recent arXiv defense benchmark that evaluates \textbf{6 novel MCP-based attack scenarios} across \textbf{6 agentic LLMs}.
\item Strong fit for your thesis because it reasons across the full lifecycle: metadata probing, runtime events, and historical traces.
\item Best used as a \textbf{benchmark design reference} and evaluation supplement rather than as a confirmed large training dataset.
\end{itemize}

\subsection*{10. Damn Vulnerable MCP Server (DVMCP)}
\textbf{File:} \texttt{20\_dvmcp.md} \\
\textbf{Source:} \url{https://github.com/harishsg993010/damn-vulnerable-MCP-server}
\begin{itemize}[leftmargin=1.2em]
\item Public practical lab with \textbf{10 MCP attack challenges} across easy, medium, and hard tiers.
\item Includes prompt injection, tool poisoning, excessive permissions, tool shadowing, token theft, malicious code execution, remote access control, and multi-vector attacks.
\item Best used for demos, red-team walk-throughs, smoke tests, and concrete defense validation against server-facing attack paths.
\end{itemize}

\subsection*{11. audit-db and vulnerability-db}
\textbf{Files:} \texttt{21\_audit\_db.md}, \texttt{22\_vulnerability\_db.md} \\
\textbf{Sources:} \url{https://github.com/ModelContextProtocol-Security/audit-db}, \url{https://github.com/ModelContextProtocol-Security/vulnerability-db}
\begin{itemize}[leftmargin=1.2em]
\item These are not classic benchmarks, but they are highly relevant supporting datasets for operational risk scoring.
\item \textbf{audit-db} defines a structured schema for reproducible MCP server audits, with manifests, findings, artifacts, and severity indexes.
\item \textbf{vulnerability-db} defines an OSV-style advisory feed for MCP vulnerabilities, which is useful for known-vulnerability priors and deployment-time enrichment.
\item Both are best treated as \textbf{external evidence feeds} rather than as standalone evaluation corpora.
\end{itemize}

\section*{If You Need Only Four Resources}
If you want the smallest high-value set for a thesis chapter or experiment section, use:
\begin{enumerate}[leftmargin=1.4em]
\item \textbf{MCPSecBench} for direct server/protocol attack evaluation.
\item \textbf{MCP-AttackBench} for large-scale attack training data.
\item \textbf{MCPTox} for realistic attacks on real MCP servers.
\item \textbf{MCP Server Database} for server capability and exfiltration-chain calibration.
\end{enumerate}

\section*{Best Combination by Goal}
\begin{longtable}{p{0.25\textwidth} p{0.63\textwidth}}
\toprule
\textbf{Goal} & \textbf{Best resources} \\
\midrule
Train malicious request detector & MCP-AttackBench, MCPSecBench \\
Evaluate direct server attack handling & MCPSecBench, MCP-SafetyBench, Component-based Attack PoC \\
Study poisoned tools on real servers & MCPTox, MCP-ITP \\
Calibrate server exposure and exfiltration risk & MCP Server Database, MCP Server Dataset (67K) \\
Support a server-defense thesis framing & MCPSecBench, Component-based Attack PoC, MCP Server Database, MCPTox \\
\bottomrule
\end{longtable}

\section*{Bottom Line}
If your main question is \textbf{attacks on the MCP server}, the strongest direct sources in this folder are:
\textbf{MCPSecBench}, \textbf{MCP-AttackBench}, \textbf{Component-based Attack PoC}, and \textbf{MCPTox}.

If your question expands to \textbf{what makes a server structurally risky even before an attack happens}, add:
\textbf{MCP Server Database} and \textbf{MCP Server Dataset (67K)}.

If you want newer or more operational internet-sourced additions, the best ones to layer on top are:
\textbf{MCPShield} for lifecycle-aware defense evaluation, \textbf{DVMCP} for practical attack labs, and
\textbf{vulnerability-db} / \textbf{audit-db} for structured external evidence.

\end{document}
